\section{Proof for three independent arms}\label{sec:spiders}

We prove Theorem~\ref{spider:stable}. The notation $Y_{a,b,c}$ and the
construction are fixed in Section~\ref{sec:bundle-families}.

\subsection{Independent arm subspaces and the central relation}

Use the independent local reflections $r\mapsto r$,
$q\mapsto-r-q$, the untwisted-ray transpositions at internal $\PP^2$
factors, and the permutations of the three central untwisted rays.
They preserve the distinguished vectors and the sums of untwisted
base coordinates, hence every twist. They also preserve adjacent
hexagon pairs and omitted base-ray sets. They therefore generate a
group $A$ of fan automorphisms for arbitrary arm lengths.

Every $A$-invariant closed ray set selects each hexagon as empty, its
distinguished line, or its whole plane, denoted $H=0,1,2$. At an
internal base the two untwisted rays are selected together, and the
twisted ray independently, with indicators $C,T$. At a leaf its two
rays are chosen independently. At the center $C_*,T_*$ select all
three untwisted rays and the twisted ray. These configurations include
every invariant flat; the arm choices need not agree.

Read each arm from the leaf to the center using the path recurrence
\eqref{path:transition}. Its state $s=G,g,u$ records whether the
outgoing distinguished line is selected, is spanned through earlier
relations, or is unspanned. Let $e$ indicate a nonempty arm selection.
Let $f=1$ when every hexagon is a plane, every internal base has $C=1$,
and the leaf has at least one ray selected. The reachable terminal
triples are
\[
 (G,1,0),\quad(G,1,1),\quad(g,1,0),\quad(u,0,0),\quad(u,1,0).
\]

\begin{lemma}\label{spider:rank}
Suppose the three arm selections have ranks $r_i$, deficits $v_i$,
and terminal triples $(s_i,e_i,f_i)$. Adding the central choice gives
rank
\[
 \sum_{i=1}^3r_i+3C_*+T_*(1-C_*)
       +C_*T_*\mathbf1_{\{\exists i:s_i=u\}}
\]
and deficit
\begin{equation}\label{spider:central-cost}
 \sum_{i=1}^3v_i+3C_*+T_*(1-C_*)
 +C_*T_*\mathbf1_{\{\exists i:s_i=u\}}
 -(3C_*+T_*)\delta_*.
\end{equation}
Here $\delta_*$ is the normalized degree of one central ray. The
selection is empty exactly when all $e_i,C_*,T_*$ vanish, and has full
rank exactly when $C_*=f_1=f_2=f_3=1$.
\end{lemma}
\begin{proof}
The arm spans occupy disjoint coordinate blocks. The central untwisted
rays add $3C_*$ dimensions. If $C_*=0$ and $T_*=1$, the twisted ray has a new
central coordinate and adds one dimension. If $C_*=T_*=1$, its remaining
image is minus the sum of the terminal distinguished vectors in the
direct sum of the three arm quotients. This sum is nonzero precisely
when at least one summand is nonzero, giving the displayed rank.
Subtracting selected weights proves~\eqref{spider:central-cost}.

Full rank requires each hexagon's nontrivial reflection representation,
each internal base's nontrivial transposition representation, and the
central base's two-dimensional standard representation. Thus every
hexagon must be a plane, every internal untwisted pair must be selected,
and $C_*=1$. The remaining leaf coordinate requires at least one leaf
ray. These conditions also suffice directly. Emptiness is immediate.
\end{proof}

Minimize each arm deficit for its five terminal triples, and combine
the results for the four central choices. This gives $5^3\cdot4=500$
entries before the empty and full-rank exclusions. Their positivity
implies stability by Proposition~\ref{prop:symmetry}. All quantitative
deficits below concern these $A$-invariant configurations.

\subsection{The mixed cubic integral}

Write $L_1=a,L_2=b,L_3=c$. On each arm number the hexagons from the
center by $p=0,\ldots,L_i-1$. With $d\nu(t)=(2-|t|)\,dt$ on $[-1,1]$
and the kernel $\mathcal K$ of~\eqref{path:kernel}, put
\[
 h_0(t)=2+t,\qquad
 h_{k+1}(s)=\int\mathcal K(s,t)h_k(t)\,d\nu(t).
\]
For positive functions $u,v$ define
\begin{align}
 J_{u,v}(s)&=\iint\frac{(4-s-t-z)^3}{6}
                u(t)v(z)\,d\nu(t)d\nu(z),\label{spider:J}\\
 J^-_{u,v}(s)&=\iint\frac{(4-s-t-z)^2}{2}
                u(t)v(z)\,d\nu(t)d\nu(z).\label{spider:Jminus}
\end{align}
Set $h^{(i)}=h_{L_i-1}$. The left initial function on arm $i$ is
$\ell_{0,i}=J_{h^{(j)},h^{(k)}}$, where $\{i,j,k\}=\{1,2,3\}$.
Propagate it by the left operator in~\eqref{path:chains}. The central
ray degree is
\begin{equation}\label{spider:central-degree}
 \delta_*=
 \frac{\int J^-_{h^{(j)},h^{(k)}}(s)h^{(i)}(s)\,d\nu(s)}
      {\int J_{h^{(j)},h^{(k)}}(s)h^{(i)}(s)\,d\nu(s)}.
\end{equation}
The value is independent of $i$. The other degrees are the path
ratios~\eqref{path:hexdegree}--\eqref{path:basedegree}, using left
index $p$ and right index $L_i-1-p$ for hexagon $p$, and left index
$c-1$ and right index $L_i-1-c$ for internal base $c$. The leaf
numerator replaces $2+t$ by $1$. These formulas follow by integrating
the other two arms in Proposition~\ref{graph:degree-formulas}.

To specify the new polynomials, put $M_r(u)=\int t^ru(t)\,d\nu(t)$.
The coefficient of $s^d$ in $J_{u,v}$ is
\begin{equation}\label{spider:coefficients}
 \sum_{\substack{r,t\ge0\\d+r+t\le3}}
 \frac{(-1)^{d+r+t}4^{3-d-r-t}M_r(u)M_t(v)}
      {d!r!t!(3-d-r-t)!}.
\end{equation}
Replacing $3$ by $2$ gives $J^-_{u,v}$. Thus the four-by-four moment
matrices already printed in Appendix~\ref{sec:paths} determine every
finite degree exactly.

\subsection{Uniform estimates with different arm lengths}

Keep $\Delta=\log(81/25)$ and $\tau=2/7$. Lemma~\ref{path:contraction}
applies to any positive initial function. The left and right chains
have positive projective limits $\ell_\infty,h_\infty$, independent
of their initial functions, with index-$k$ error at most
$\Delta\tau^{k-1}$ for $k\ge1$. Each degree ratio has logarithmic
error at most the sum of the projective errors of its factors.

Every individual ray weight in the following actual and comparison
configurations is at most $B=81$. Indeed, for fixed $t,z$ the
oscillation ratio of $(4-s-t-z)^3$ is at most $27$, and this persists
under positive integration. Right functions have oscillation ratio
at most $9$. Hexagon facet measures have mass one and $\nu$ has mass
three, giving $27\cdot9/3=81$. Internal, leaf and central base weights
are positive averages bounded by $2$, $1$ and $3$ respectively.

Fix $q=16$. Call an arm short if $2\le L_i<2q$, and otherwise write
$L_i=2q+m_i$, $m_i\ge0$. First replace the incoming function $h^{(i)}$
of every long arm by $h_\infty$, keeping short incoming functions
unchanged. Denote these functions by $\widehat h^{(i)}$. Use them in
the central degree and in each initial left function; for the moment
retain every actual right chain on the arms.

Put $\bar\Delta=11757/10000>\Delta$, using the rational exponential
estimate of Appendix~\ref{sec:paths}. For a long arm set
$\epsilon_i=\bar\Delta\tau^{L_i-2}$, and for a short arm set
$\epsilon_i=0$. Positivity of the mixed integral gives
\[
 d_H\bigl(J_{h^{(j)},h^{(k)}},
          J_{\widehat h^{(j)},\widehat h^{(k)}}\bigr)
 \le\epsilon_j+\epsilon_k.
\]
Let $z=3\bar\Delta\tau^{2q-2}$. Each sum on the right is at most $z$.
After propagation along arm $i$, the logarithmic weight change is at
most $z\tau^p$ for hexagon $p$, $z\tau^{c-1}$ for internal base $c$,
and $z\tau^{L_i-1}$ for the leaf. The central logarithmic change is
at most $\sum_i\epsilon_i\le z$.

Using $e^t-1\le t+t^2$ for $0\le t\le1$, sum the six hexagon rays,
three internal-base rays and two leaf rays on each arm. Their geometric
sum is bounded by $9/(1-\tau)+2$. The four central rays give the
uniform junction error
\begin{equation}\label{spider:junction-error}
 E_J=B\left(10+\frac{27}{1-\tau}\right)(z+z^2).
\end{equation}
In particular, increasing the length of another arm does not multiply
this perturbation by its number of rays.

Next replace each long arm by three zones. Read an internal period as
base $c$ followed by hexagon $c-1$ and give both the index $p=c-1$;
hexagon $p$ has the same zone index. For $p<q$, keep the finite left
function and use $h_\infty$ on the right. For $q\le p<q+m_i$, use
both limiting functions. For $p\ge q+m_i$, keep the right function
and use $\ell_\infty$ on the left. The leaf uses the limiting left
function. Short arms retain the preceding mixed-junction weights.
Write $M(m_1,\ldots,m_k)$ for this comparison, where $k$ is the number
of long arms. Its initial junction functions depend on the short
lengths and the set of long arms, but not on any $m_i$.

Put $x=\bar\Delta\tau^{q-2}$ and $S=2x/(1-\tau)$. Each boundary zone
has at most $9q+4$ rays and omitted-chain logarithmic error at most $x$.
In a middle period $p$ the error is at most
\[
 t_p=\bar\Delta(\tau^{p-1}+\tau^{L_i-3-p}),\qquad
 t_p\le2x,\qquad \sum_pt_p\le S.
\]
There are at most nine rays per period. Summing over at most three
long arms gives
\begin{equation}\label{spider:zone-error}
 E_Z=3B\{9S(1+2x)+2(9q+4)(x+x^2)\}.
\end{equation}
Hence every selected-ray deficit changes by at most $E_J+E_Z$ from
the actual geometry to its comparison, uniformly in all arm lengths.

\subsection{Finite comparisons and independent arm shortening}

Consider the thirty short lengths $2,\ldots,31$ and the four long-arm
middle lengths $m=0,1,2,3$. Their unordered triples give $7140$ finite
comparisons. Ordering removes duplicate geometries only; every case
still includes independent selections on its arms.

For rational evaluation put $P=40$. Replace $h_\infty$ by $h_P$ and
$\ell_\infty$ by the index-$P$ left chain starting at $2-s$. The
common limit is independent of this initial choice. With
$y=3\bar\Delta\tau^{P-1}$, each ray degree has logarithmic error at
most $y$: a finite mixed left function changes by at most
$2\bar\Delta\tau^{P-1}$, and an omitted right or left limit by at
most $\bar\Delta\tau^{P-1}$. A finite comparison has at most
$27(2q+3)+1$ rays, so its deficit error is bounded by
\begin{equation}\label{spider:proxy-error}
 E_R=B\{27(2q+3)+1\}(y+y^2).
\end{equation}

Let $\mu$ be the smallest proper nonempty rational deficit in these
comparisons, and let $F_0$ be the largest absolute all-selected cost.
Table~\ref{spider:table} records the finite minima by number of long
arms; each displayed lower bound is rounded down. The arm recurrence,
the five terminal states, \eqref{spider:central-cost}, and the mixed
moment formula specify the finite calculation completely.

% The reviewed finite table and bounds are inserted after computation.
\begin{table}[ht]
\centering
\caption{Strict rational lower bounds for proper nonempty configurations
in the finite mixed-arm comparisons. All short lengths are $2,\ldots,31$;
each long arm has middle length $0,1,2,3$.}
\label{spider:table}
\begin{tabular}{rrr}
\toprule
Long arms & Comparisons & Lower bound \\
\midrule
0 & 4960 & 0.029198008 \\
1 & 1860 & 0.029198008 \\
2 & 300 & 0.029198008 \\
3 & 20 & 0.029505518 \\
\bottomrule
\end{tabular}
\end{table}
The exact rational calculations give
\[
 \mu>0.029198,\quad \gamma>0.1259,\quad F_0<10^{-15},
\]
and
\[
 E_J<6.506\cdot10^{-13},\qquad
 E_Z<0.002216796,\qquad E_R<1.634\cdot10^{-16}.
\]


At true limiting weights, the empty $u$-loop and all-selected $G$-loop
have cost zero. The latter assertion follows because a limiting period
has total weight four, as proved in Appendix~\ref{sec:paths}. The same
period calculation gives a positive rational simple-cycle gap $\gamma$,
after the two zero loops are removed. A simple cycle uses at most
three nine-ray periods, so its true gap is at least $\gamma-E_R$.

Set $F=F_0+E_R$. Inserting a full middle period changes neither the
all-selected cost nor the junction data of a true comparison. Therefore
the all-selected cost is at least $-F$ for every choice of middle
lengths. Every full-rank selection has cost at least this value:
adding its omitted positive-weight rays changes no rank. Empty
selections have cost zero.

\begin{proof}[Proof of Theorem~\ref{spider:stable}]
We first prove that every proper nonempty true comparison has deficit
at least
\[
 \min\{\mu-E_R,\ \gamma-E_R-F\}.
\]
Induct on the sum of its middle lengths. If each is at most three,
the finite comparison applies. Otherwise choose an arm with at least
four middle periods. Two of its boundary states agree among $G,g,u$.
Delete the closed walk between them on this arm. The retained indices
from the center and leaf are unchanged, as are the junction functions,
so the resulting weights are those of the shorter comparison.

If the deleted walk contains only the empty $u$-loop, it removes no
selected ray, and the retained $u$ boundary prevents full rank. If it
contains only the full $G$-loop, the retained $G$ boundary prevents
emptiness, and deleting selected planes and untwisted pairs removes
no failure of a full-rank condition. Proper nonemptiness is therefore
preserved in both zero-cost cases, and induction applies.

Every other closed walk contains a positive simple cycle, with cost
at least $\gamma-E_R$. The remaining configuration is proper nonempty,
empty, or full rank. Apply induction, the zero empty cost, or the bound
$-F$, respectively. This proves the comparison inequality; it uses
the gap of one arm, without an equal-length or simultaneous-deletion
assumption.

Subtracting the actual-to-comparison error gives every proper nonempty
$A$-invariant configuration deficit at least
\[
 \min\{\mu-E_R,\ \gamma-E_R-F\}-E_J-E_Z>0.02698.
\]
The symmetry criterion proves anticanonical stability. The dimension
and Picard number follow from the graph construction.
\end{proof}
